%% toc: P15 174 174.5,174.6,174.7,174.8
\section{Proofs for treatment versions, interference, adherence, missing outcomes and sample size}
\label{sec:annex-p15-w02b-p15-076}

Notation is that of \cref{def:p15-w02b-p15-076-setting} and of the
definition attached to each result. All random variables in the first four
proofs are functions on a finite outcome space. For an event \(E\) of
positive probability and a real random variable \(X\),
\(\mathbb E[X\mid E]=\mathbb E[X\mathbf 1_E]/\Pr(E)\). If random variables
\(A\) and \(B\) are independent and \(\alpha,\beta\) are real functions,
then \(\mathbb E[\alpha(A)\beta(B)]=\mathbb E[\alpha(A)]\,\mathbb E[\beta(B)]\);
and a function of a tuple independent of \(A\) is independent of \(A\).
Both facts are proved as (F2) and (F3) in
\cref{sec:annex-p15-w02b-p15-077-assignment-scm}.

\begin{annexitem}{annex:p15-w02b-p15-076-weighting}{thm:p15-w02b-p15-076-weighting}{The weighting statistic under any observation model}
\begin{proof}
For each unit \(i\) and cell \(t\), \(\Pr(Z_i=t)=\pi_i(t)>0\), so
\(\mathbb E[\mathbf 1_{\{Z_i=t\}}Y_i^{\mathrm{rec}}]
 =\pi_i(t)\,\mathbb E[Y_i^{\mathrm{rec}}\mid Z_i=t]\).
Linearity of expectation over the finite double sum defining
\(\widehat\tau\) gives
\[
 \mathbb E[\widehat\tau]=\frac1N\sum_i\sum_t
  \frac{w_t\,\pi_i(t)\,\mathbb E[Y_i^{\mathrm{rec}}\mid Z_i=t]}{\pi_i(t)},
\]
which is the stated formula. If
\(\mathbb E[Y_i^{\mathrm{rec}}\mid Z_i=t]=Y_i(t)\) whenever \(w_t\ne0\),
the right side is \(N^{-1}\sum_i\sum_tw_tY_i(t)=\tau\); terms with
\(w_t=0\) contribute zero to both sides.
\end{proof}
\end{annexitem}

\begin{annexitem}{annex:p15-w02b-p15-076-versions}{thm:p15-w02b-p15-076-versions}{Version mixtures}
\begin{proof}
(1) On the event \(\{Z_i=t\}\),
\(Y_i^{\mathrm{rec}}=Y_i(t,\phi_i(t,\xi_i))\). The tuple \(\xi\) is
independent of \(Z\), so \(\xi_i\) is independent of \(Z_i\), and
\[
 \mathbb E\bigl[\mathbf 1_{\{Z_i=t\}}Y_i^{\mathrm{rec}}\bigr]
 =\mathbb E\bigl[\mathbf 1_{\{Z_i=t\}}\bigr]\,
  \mathbb E\bigl[Y_i(t,\phi_i(t,\xi_i))\bigr]
 =\pi_i(t)\sum_{v\in\mathcal V_t}\lambda_i(v\mid t)Y_i(t,v),
\]
the last step grouping outcomes by the value \(v\) of
\(\phi_i(t,\xi_i)\). Hence
\(\mathbb E[Y_i^{\mathrm{rec}}\mid Z_i=t]=\sum_v\lambda_i(v\mid t)Y_i(t,v)\),
and \cref{thm:p15-w02b-p15-076-weighting} gives
\(\mathbb E[\widehat\tau]=\tau_\lambda\).

(2) If \(Y_i(t,v)=Y_i(t)\) for every \(v\in\mathcal V_t\), then
\(\sum_v\lambda_i(v\mid t)Y_i(t,v)=Y_i(t)\sum_v\lambda_i(v\mid t)=Y_i(t)\),
because \(\lambda_i(\cdot\mid t)\) is a probability mass function on
\(\mathcal V_t\). Thus \(\tau_\lambda=\tau\). Also
\(Y_i(t,v_t)=Y_i(t)\) for every selection, so \(\theta_{\mathbf v}=\tau\).

(3) Choose \(v_+,v_-\in\mathcal V_t\) with
\(Y_i(t,v_+)>Y_i(t,v_-)\). Let \(\xi\) be constant, which is independent of
every random variable. For every unit \(j\) and cell \(s\) fix a version
\(f_j(s)\in\mathcal V_s\), and define the first mechanism by
\(\phi_j(s,\cdot)=f_j(s)\) for \((j,s)\ne(i,t)\) and
\(\phi_i(t,\cdot)=v_+\); define the second identically except that
\(\phi_i(t,\cdot)=v_-\). The version laws are point masses, and the two
values of \(\tau_\lambda\) differ only in the term of unit \(i\) and cell
\(t\):
\[
 \tau_\lambda-\tau_{\lambda'}=\frac{w_t}{N}\bigl(Y_i(t,v_+)-Y_i(t,v_-)\bigr)\ne0.
\]

(4) Let \(\sigma\) be the permutation of \(\mathcal V_t\) exchanging
\(v_t\) and \(v'\). The second model keeps the assignment law and \(\xi\),
sets \(Y_i'(t,v)=Y_i(t,\sigma(v))\) and
\(\phi_i'(t,x)=\sigma(\phi_i(t,x))\), and leaves every other potential
response and map unchanged. At an outcome with \(Z_i=t\),
\[
 Y_i^{\prime\,\mathrm{rec}}
 =Y_i'\bigl(t,\sigma(\phi_i(t,\xi_i))\bigr)
 =Y_i\bigl(t,\sigma(\sigma(\phi_i(t,\xi_i)))\bigr)
 =Y_i\bigl(t,\phi_i(t,\xi_i)\bigr)=Y_i^{\mathrm{rec}},
\]
since \(\sigma\circ\sigma\) is the identity. At an outcome with
\(Z_i\ne t\), and for every other unit, nothing changed. So the recorded
vector \((Z,Y^{\mathrm{rec}})\) is the same at every outcome. The target
changes only through unit \(i\) and cell \(t\):
\[
 \theta'_{\mathbf v}-\theta_{\mathbf v}
 =\frac{w_t}{N}\bigl(Y_i(t,\sigma(v_t))-Y_i(t,v_t)\bigr)
 =\frac{w_t}{N}\bigl(Y_i(t,v')-Y_i(t,v_t)\bigr)\ne0.
\]
The new version law is
\(\lambda_i'(v\mid t)=\Pr(\sigma(\phi_i(t,\xi_i))=v)=\lambda_i(\sigma(v)\mid t)\);
it agrees with \(\lambda_i(\cdot\mid t)\) at every version other than
\(v_t,v'\), and at those two exactly when
\(\lambda_i(v_t\mid t)=\lambda_i(v'\mid t)\).

(5) On \(\{Z_i=t,V_i=v\}\) the recorded response is \(Y_i(t,v)\), and
\(\{Z_i=t,V_i=v\}=\{Z_i=t\}\cap\{\phi_i(t,\xi_i)=v\}\). By independence of
\(Z_i\) and \(\xi_i\),
\[
 \begin{aligned}
 \Pr(Z_i=t,V_i=v)&=\pi_i(t)\lambda_i(v\mid t),\\
 \mathbb E\bigl[\mathbf 1_{\{Z_i=t,V_i=v\}}Y_i^{\mathrm{rec}}\bigr]
  &=\pi_i(t)\lambda_i(v\mid t)Y_i(t,v).
 \end{aligned}
\]
When \(\lambda_i(v_t\mid t)>0\), the ratio in the statement equals
\(Y_i(t,v_t)\); summing with weights gives \(\theta_{\mathbf v}\), and
linearity shows that the statistic with realized indicators has
expectation \(\theta_{\mathbf v}\). Its denominators
\(\pi_i(t)\lambda_i(v_t\mid t)\) are probabilities of recorded events, so
they are fixed by the joint law of \((Z,V)\). If
\(\lambda_i(v_t\mid t)=0\), the value \(Y_i(t,v_t)\) enters the recorded
variables only on the event \(\{Z_i=t,V_i=v_t\}\), of probability
\(\pi_i(t)\cdot0=0\). Changing that value therefore changes the recorded
variables at most on a null set and leaves their joint law unchanged, while
\(\theta_{\mathbf v}\) changes by \(w_t/N\) times the change.
\end{proof}
\end{annexitem}

\begin{annexitem}{annex:p15-w02b-p15-076-interference}{thm:p15-w02b-p15-076-interference}{Design averages}
\begin{proof}
(1) Grouping outcomes by the value \(z\) of \(Z\),
\[
 \mathbb E\bigl[\mathbf 1_{\{Z_i=t\}}Y_i(Z)\bigr]
 =\sum_{z:\,z_i=t}\varpi(z)Y_i(z).
\]
Dividing by \(\pi_i(t)\) gives \(\mathbb E[Y_i(Z)\mid Z_i=t]\), and
\cref{thm:p15-w02b-p15-076-weighting} gives
\(\mathbb E[\widehat\tau]=\tau_\varpi\).

(2) Under no interference \(Y_i(z)=Y_i(t)\) whenever \(z_i=t\), so
\[
 \sum_{z:\,z_i=t}\varpi(z)Y_i(z)=Y_i(t)\sum_{z:\,z_i=t}\varpi(z)
 =Y_i(t)\pi_i(t).
\]
Hence \(\tau_\varpi=\tau\). Also \(Y_i(t\mathbf 1)=Y_i(t)\), because the
\(i\)th coordinate of \(t\mathbf 1\) is \(t\); so \(\tau_{\rm com}=\tau\).

(3) Take \(N=2\), \(\mathcal T=\{0,1\}\), \(w=(-1,1)\) and
\(Y_1(z)=Y_2(z)=z_1+z_2\). Then \(\tau_{\rm com}=\tfrac12\bigl((2-0)+(2-0)\bigr)=2\).
Under the law \(\varpi_A\) with mass \(1/2\) on each of \((1,0)\) and
\((0,1)\), \(\pi_i(1)=\pi_i(0)=1/2\). For unit \(1\), the only vector with
\(z_1=1\) is \((1,0)\), with response \(1\), and the only vector with
\(z_1=0\) is \((0,1)\), with response \(1\); unit \(2\) is symmetric. So
\(\tau_{\varpi_A}=0\). Under the law \(\varpi_B\) with mass \(1/4\) on each
of the four vectors, again \(\pi_i(t)=1/2\). For unit \(1\),
\(\mathbb E[Y_1(Z)\mid Z_1=1]=\tfrac12(1+2)=3/2\) and
\(\mathbb E[Y_1(Z)\mid Z_1=0]=\tfrac12(0+1)=1/2\); unit \(2\) is
symmetric. So \(\tau_{\varpi_B}=1\). The three values \(0,1,2\) are
distinct, and the two designs have the same marginals.

(4) Let \(\varpi(z)>0\) and \(z_i=t\). Because the design assigns by
cluster, \(z_j=t\) for every \(j\in g(i)\), so \(z\) and \(t\mathbf 1\)
agree on \(g(i)\); confinement of interference gives
\(Y_i(z)=Y_i(t\mathbf 1)\). Vectors of probability zero do not contribute
to the sum, so
\[
 \sum_{z:z_i=t}\varpi(z)Y_i(z)=Y_i(t\mathbf 1)\sum_{z:z_i=t}\varpi(z)
 =Y_i(t\mathbf 1)\pi_i(t),
\]
and \(\tau_\varpi=N^{-1}\sum_i\sum_tw_tY_i(t\mathbf 1)=\tau_{\rm com}\).
By part~1 the weighting statistic is unbiased for \(\tau_{\rm com}\).
\end{proof}
\end{annexitem}

\begin{annexitem}{annex:p15-w02b-p15-076-adherence}{thm:p15-w02b-p15-076-adherence}{The complier identity and its two premises}
\begin{proof}
(1) On \(\{Z_i=t\}\), \(D_i=D_i(t)\) and
\(Y_i^{\mathrm{rec}}=Y_i(t,D_i(t))\), both constants. Hence
\(\mathbb E[D_i\mid Z_i=t]=D_i(t)\) and
\(\mathbb E[Y_i^{\mathrm{rec}}\mid Z_i=t]=Y_i(t,D_i(t))\).
\Cref{thm:p15-w02b-p15-076-weighting} with \(w=(-1,1)\) gives the
expectations \(\mathrm{ITT}_D\) and \(\mathrm{ITT}_Y\). No property of
\(D_i(\cdot)\) or \(Y_i(\cdot,\cdot)\) was used.

(2) Under exclusion the unit term of \(\mathrm{ITT}_Y\) is
\(Y_i(D_i(1))-Y_i(D_i(0))\). Check the four types. If
\(D_i(1)=D_i(0)\), the term and \((D_i(1)-D_i(0))(Y_i(1)-Y_i(0))\) are
both zero. For a complier the term is \(Y_i(1)-Y_i(0)\) and
\(D_i(1)-D_i(0)=1\). For a defier the term is \(Y_i(0)-Y_i(1)\) and
\(D_i(1)-D_i(0)=-1\). Summing over units gives the identity.

(3) Monotonicity excludes defiers, so \(D_i(1)-D_i(0)\) is one for a
complier and zero otherwise. Hence \(\mathrm{ITT}_D=|\mathcal C|/N\), and by
part~2 \(\mathrm{ITT}_Y=N^{-1}\sum_{i\in\mathcal C}(Y_i(1)-Y_i(0))\). If
\(\mathcal C\ne\varnothing\), the last sum is
\(|\mathcal C|\cdot\mathrm{CACE}\), so
\(\mathrm{ITT}_Y=\mathrm{ITT}_D\cdot\mathrm{CACE}\); if
\(\mathcal C=\varnothing\), both contrasts are zero. When
\(\mathrm{ITT}_D>0\) the set \(\mathcal C\) is nonempty and division gives
the ratio.

(4) If \(D_i(0)=0\), then \(D_i(1)\geq0=D_i(0)\), so monotonicity holds;
a unit is a complier exactly when \(D_i(1)=1\). Part~3 applies.

(5) Take two units. Unit \(1\) has \(D_1(0)=0\), \(D_1(1)=1\) and
\(Y_1(z,d)=0\) for all \(z,d\): a complier whose response is unaffected by
assignment and receipt. Unit \(2\) has \(D_2(0)=D_2(1)=0\),
\(Y_2(1,0)=1\), \(Y_2(0,0)=0\) and \(Y_2(z,1)=0\): a never-receiver for
whom exclusion fails. Monotonicity holds,
\(\mathrm{ITT}_D=(1+0)/2=1/2\), and
\(\mathrm{ITT}_Y=\bigl((0-0)+(1-0)\bigr)/2=1/2\), so the ratio is \(1\).

(6) Take three units with exclusion. Units \(1\) and \(2\) are compliers
with \(Y_i(0)=0\) and \(Y_i(1)=2\); unit \(3\) is a defier with
\(Y_3(0)=0\) and \(Y_3(1)=4\). Every unit has a positive receipt effect.
By part~2, \(\mathrm{ITT}_Y=(2+2-4)/3=0\), while
\(\mathrm{ITT}_D=(1+1-1)/3=1/3\) and \(\mathrm{CACE}=2\).
\end{proof}
\end{annexitem}

\begin{annexitem}{annex:p15-w02b-p15-076-missing}{thm:p15-w02b-p15-076-missing}{Bounds, bias and strata}
\begin{proof}
(1) On \(\{Z_i=t\}\), \(R_i=R_i(t)\) and
\(R_iY_i(Z_i)=R_i(t)Y_i(t)\), both constants. Hence
\(\mathbb E[\mathbf 1_{\{Z_i=t\}}R_iY_i(Z_i)]/\pi_i(t)=R_i(t)Y_i(t)\) and
\(\mathbb E[\mathbf 1_{\{Z_i=t\}}R_i]/\pi_i(t)=R_i(t)\); averaging over
units gives \(A_t\) and \(\bar r_t\).

(2) Split the sum defining \(\mu(t)\):
\[
 \mu(t)=A_t+\frac1N\sum_{i:R_i(t)=0}Y_i(t).
\]
The second sum has \(Nq_t\) terms, each an unrecorded response free to take
any value in \([\ell,u]\). Each term lies in \([\ell,u]\), so the sum lies
in \([Nq_t\ell,Nq_tu]\). Conversely, setting every unrecorded response
equal to a common \(y\in[\ell,u]\) gives the sum \(Nq_ty\), which runs over
the whole interval as \(y\) does. This proves the first claim, including
the degenerate case \(q_t=0\). For the contrast,
\[
 \tau=\sum_tw_tA_t+\sum_t\frac{w_t}{N}\sum_{i:R_i(t)=0}Y_i(t).
\]
Distinct cells involve distinct unrecorded responses, so the inner sums
vary independently. By the first claim, the term of cell \(t\) ranges over
\(\{w_tq_ty:y\in[\ell,u]\}\), the interval from
\(q_t\min(w_t\ell,w_tu)\) to \(q_t\max(w_t\ell,w_tu)\). A sum of
independently varying terms, each ranging over a closed interval, ranges
over the interval whose endpoints are the sums of the endpoints: the
extreme values are attained together, and every intermediate value is
attained by moving one term at a time continuously. Each cell contributes
the width \(q_t|w_t|(u-\ell)\).

(3) Since \(\operatorname{Cov}_N(R(t),Y(t))=A_t-\bar r_t\mu(t)\) and
\(\mu_{\rm cc}(t)=A_t/\bar r_t\),
\[
 \mu_{\rm cc}(t)-\mu(t)=\frac{A_t-\bar r_t\mu(t)}{\bar r_t}
 =\frac{\operatorname{Cov}_N(R(t),Y(t))}{\bar r_t}.
\]
If \(q_t>0\), splitting \(N\mu(t)\) into recorded and unrecorded parts gives
\(\mu(t)=\bar r_t\mu_{\rm cc}(t)+q_t\mu_{\rm mis}(t)\), so
\(\mu_{\rm cc}(t)-\mu(t)=(1-\bar r_t)\mu_{\rm cc}(t)-q_t\mu_{\rm mis}(t)
=q_t(\mu_{\rm cc}(t)-\mu_{\rm mis}(t))\). Both means are averages of values
in \([\ell,u]\), so their difference has absolute value at most
\(u-\ell\). The first display vanishes exactly when the covariance does.

(4) Apply part~3 within stratum \(x\): the hypothesis gives
\(\mu_{\rm cc}(t,x)=\mu(t,x)\), the mean of \(Y_i(t)\) over the stratum.
Grouping the units of \(\mu(t)\) by stratum,
\(\mu(t)=\sum_x(N_x/N)\mu(t,x)=\sum_x(N_x/N)\mu_{\rm cc}(t,x)\).
For the final claim take ten units in cell \(t\) with range \([0,4]\):
five low-stratum units with responses \(1,2,3,2,2\), the first three
recorded; and five high-stratum units with responses \(0,0,1,1,1\), all
recorded. In the low stratum the recorded and unrecorded means are both
\(2\), so the stratum covariance is zero by part~3; in the high stratum
every unit is recorded, so the covariance is zero. The premise holds, and
\(\mu(t)=13/10=1.3\), whereas
\(\mu_{\rm cc}(t)=(1+2+3+0+0+1+1+1)/8=9/8\).
\end{proof}
\end{annexitem}

The remaining proofs use the normal linear model of
\cref{thm:p15-factorial-linear-uncertainty}. Its proof, in
\cref{annex:p15.factorial.linear-uncertainty}, chooses an orthonormal
basis whose first vector is \(a_*/\|a_*\|\) and whose last \(\nu\) vectors
span the residual space. The coordinates of \(\epsilon/\sigma\) in that
basis are independent standard normal variables. Hence
\[
 Z_c=\frac{\widehat\theta-\theta}{\sigma\|a_*\|}\sim N(0,1)
 \quad\text{and}\quad
 W=\frac{\nu s^2}{\sigma^2}=\frac{\mathrm{SSE}}{\sigma^2}\sim\chi^2_\nu
\]
are independent, and \(\mathbb E[s^2]=\sigma^2\). For a function
\(F\) bounded and measurable on \(\mathbb R\times(0,\infty)\), independence
and Fubini's theorem for the product density give
\(\mathbb E[F(Z_c,W)]=\int_0^\infty\mathbb E[F(Z_c,w)]\,dG_\nu(w)\).
Since \(W\) has a density, \(\Pr(W=0)=0\), and so \(s>0\) with probability
one.

\begin{annexitem}{annex:p15-w02b-p15-076-families}{thm:p15-w02b-p15-076-families}{Three planning families}
\begin{proof}
(1) The row space of \(X_{(n)}\) is that of \(X_1\), so \(c\) is estimable
for \(X_{(n)}\) by \cref{thm:p15-factorial-estimability}, and
\(\operatorname{rank}X_{(n)}=r\). Write \(a_{*,1}=X_1b\). The vector
\(a=n^{-1}(a_{*,1},\ldots,a_{*,1})\) equals \(X_{(n)}b/n\), so it lies in
\(\operatorname{col}(X_{(n)})\), and
\(X_{(n)}^{\mathsf T}a=\sum_{j=1}^nX_1^{\mathsf T}a_{*,1}/n=c\). By the
uniqueness in \cref{thm:p15-factorial-linear-uncertainty}, \(a=a_{*,n}\),
and \(\|a\|^2=n\|a_{*,1}\|^2/n^2\). The residual degrees of freedom are
\(nN_1-r\). For the saturated full factorial, \(X_1=H\) has
\(H^{\mathsf T}H=2^kI\) by \cref{thm:p15-factorial-walsh}, so for
\(c=e_S\), \(a_{*,1}=H(H^{\mathsf T}H)^{-1}e_S=2^{-k}\chi_S\), with
\(\|a_{*,1}\|^2=2^{-2k}\cdot2^k=2^{-k}\); and \(r=2^k\), so
\(\nu_n=n2^k-2^k\).

(2) By \cref{thm:p15-factorial-blocks}, the estimate is
\(\sum_tc_t\overline Y_{t\cdot}=\sum_{t,b}(c_t/B)Y_{tb}\), with variance
\(\sigma^2\sum_tc_t^2/B\) and \((L-1)(B-1)\) residual degrees of freedom.
Its coefficient vector depends on the treatment only, so it is a
combination of the treatment indicator columns and lies in the column
space; it is therefore \(a_{*,B}\), with
\(\|a_{*,B}\|^2=B\sum_tc_t^2/B^2=\sum_tc_t^2/B\).

(3) The cluster mean is
\(\overline Y_g=\mu_{t(g)}+U_g+\overline e_g\), with
\(\overline e_g=m^{-1}\sum_je_{gj}\). The independent terms \(U_g\) and
\(\overline e_g\) have variances \(\sigma_U^2\) and \(\sigma_e^2/m\), so
\[
 \operatorname{Var}(\overline Y_g)=\sigma_U^2+\frac{\sigma_e^2}{m}
 =\sigma^2\Bigl(\rho+\frac{1-\rho}m\Bigr)
 =\frac{\sigma^2(1+(m-1)\rho)}m=\sigma_c^2.
\]
Distinct cluster means are independent normal variables, since they are
functions of disjoint sets of independent normal terms. Thus the \(2K\)
cluster means follow the normal linear model with two cell-indicator
columns and error variance \(\sigma_c^2\). For the contrast with weights
\((-1,1)\), the least-squares vector has entries \(\pm1/K\), norm squared
\(2/K\), so the difference of cell means has variance
\(2\sigma_c^2/K=2\sigma^2(1+(m-1)\rho)/(mK)\); the model has rank \(2\),
leaving \(2K-2\) residual degrees of freedom. With individual assignment
of \(Km\) units per cell and independent errors of variance
\(\sigma^2\), each cell mean has variance \(\sigma^2/(Km)\) and the
difference \(2\sigma^2/(Km)\).
\end{proof}
\end{annexitem}

\begin{annexitem}{annex:p15-w02b-p15-076-known-variance}{thm:p15-w02b-p15-076-known-variance}{Normal precision and power with a known variance}
\begin{proof}
Write \(\omega_n=\sigma\sqrt{\kappa_1/n}\). By the preliminary paragraph,
\(\widehat\theta_n=\theta+\omega_nZ_c\) with \(Z_c\sim N(0,1)\).

(1) \(\Pr(|\widehat\theta_n-\theta|\leq z_{1-\eta/2}\omega_n)
=\Phi(z_{1-\eta/2})-\Phi(-z_{1-\eta/2})=1-\eta\). The half-width
\(z_{1-\eta/2}\omega_n\) is at most \(h\) exactly when
\(n\geq z_{1-\eta/2}^2\sigma^2\kappa_1/h^2\), since both sides are
positive and \(\omega_n\) decreases in \(n\).

(2) At a value \(\theta\),
\((\widehat\theta_n-\theta_0)/\omega_n=Z_c+(\theta-\theta_0)/\omega_n\), so
the rejection probability is
\(\Pr(Z_c>z_{1-\eta}-(\theta-\theta_0)/\omega_n)
=\Phi((\theta-\theta_0)/\omega_n-z_{1-\eta})\). At
\(\theta=\theta_0+\delta\) this is \(\Phi(d_n-z_{1-\eta})\). For
\(\theta\leq\theta_0\) it is at most \(\Phi(-z_{1-\eta})=\eta\), because
\(\Phi\) is increasing. Because \(\Phi\) is continuous and strictly
increasing, \(\Phi(d_n-z_{1-\eta})\geq1-\gamma\) exactly when
\(d_n\geq z_{1-\eta}+z_{1-\gamma}\), that is
\(\sqrt n\geq(z_{1-\eta}+z_{1-\gamma})\sigma\sqrt{\kappa_1}/\delta\);
both sides are positive, and squaring gives the bound.

(3) At \(\theta_0+\delta\) the statistic is \(Z_c+d_n\), so
\[
 \Pr(|Z_c+d_n|>z_{1-\eta/2})
 =\Phi(d_n-z_{1-\eta/2})+\Phi(-d_n-z_{1-\eta/2}),
\]
which is symmetric in \(d_n\) and so may be written with \(|d_n|\). At
\(d_n=0\) it is \(2\Phi(-z_{1-\eta/2})=\eta\). The second term is
nonnegative, so the power is at least \(\Phi(|d_n|-z_{1-\eta/2})\), and
the argument of part~2 with \(z_{1-\eta/2}\) in place of \(z_{1-\eta}\)
gives the sufficient size.
\end{proof}
\end{annexitem}

\begin{annexitem}{annex:p15-w02b-p15-076-student}{thm:p15-w02b-p15-076-student}{Student planning}
\begin{proof}
(1) Let \(Z_0\sim N(0,1)\) and \(W_0\sim\chi^2_\nu\) be independent. By
\cref{thm:p15-student-law}, \(Z_0/\sqrt{W_0/\nu}\) has distribution
function \(F_\nu\). For \(z>0\), the preliminary paragraph applied to the
indicator of \(\{Z_0\leq z\sqrt{W_0/\nu}\}\) gives
\[
 F_\nu(z)=\int_0^\infty\Phi\bigl(z\sqrt{w/\nu}\bigr)\,dG_\nu(w)
        =\mathbb E\bigl[\psi(W_0/\nu)\bigr],\qquad
 \psi(x)=\Phi(z\sqrt x).
\]
On \((0,\infty)\), \(\psi'(x)=z\varphi(z\sqrt x)/(2\sqrt x)\), where
\(\varphi\) is the standard normal density. Both factors
\(\varphi(z\sqrt x)\) and \(1/\sqrt x\) are positive and strictly
decreasing in \(x\), so \(\psi'\) is strictly decreasing and \(\psi\) is
strictly concave. Hence \(\psi(x)\leq\psi(1)+\psi'(1)(x-1)\), with strict
inequality for \(x\ne1\). The variable \(X=W_0/\nu\) has mean \(1\),
because \(W_0\) is a sum of \(\nu\) squares of standard normal variables,
and it has a density, so \(\Pr(X\ne1)=1\). Taking expectations gives
\(F_\nu(z)=\mathbb E\psi(X)<\psi(1)=\Phi(z)\). For \(1/2<p<1\),
\(z_p>0\), so \(F_\nu(z_p)<p=F_\nu(t_{\nu,p})\); as \(F_\nu\) is strictly
increasing, \(z_p<t_{\nu,p}\). The pivot
\((\widehat\theta_n-\theta)/(s_n\sqrt{\kappa_1/n})=Z_c/\sqrt{W/\nu}\) has
law \(F_\nu\), so with \(z=z_{1-\eta/2}\) and the symmetry
\(F_\nu(-z)=1-F_\nu(z)\), the coverage is
\(F_\nu(z)-F_\nu(-z)=2F_\nu(z)-1<2\Phi(z)-1=1-\eta\).

(2) By the Cauchy--Schwarz inequality,
\(\mathbb E[s_n]\leq(\mathbb E[s_n^2])^{1/2}=\sigma\), so
\(\mathbb E[H_n]=t_{\nu,1-\eta/2}\sqrt{\kappa_1/n}\,\mathbb E[s_n]
\leq t_{\nu,1-\eta/2}\sigma\sqrt{\kappa_1/n}\). Next,
\(H_n\leq h\) exactly when
\(s_n^2\leq nh^2/(t_{\nu,1-\eta/2}^2\kappa_1)\), that is, when
\(W=\nu s_n^2/\sigma^2\leq\nu nh^2/(t_{\nu,1-\eta/2}^2\sigma^2\kappa_1)\).
Since \(W\sim\chi^2_\nu\), the probability is \(G_\nu\) at that point.

(3) At \(\theta=\theta_0+\delta\),
\(\widehat\theta_n-\theta_0=\delta+\sigma\sqrt{\kappa_1/n}\,Z_c\) and
\(s_n=\sigma\sqrt{W/\nu}\), so on \(\{W>0\}\),
\[
 T=\frac{Z_c+d_n}{\sqrt{W/\nu}},\qquad
 \begin{aligned}
 \{T>q\}&=\{Z_c>q\sqrt{W/\nu}-d_n\},\\
 \{T<-q\}&=\{Z_c<-q\sqrt{W/\nu}-d_n\}.
 \end{aligned}
\]
The preliminary paragraph, applied to the indicators of these events,
gives the two integrals, since
\(\Pr(Z_c>q\sqrt{w/\nu}-d_n)=\Phi(d_n-q\sqrt{w/\nu})\) and
\(\Pr(Z_c<-q\sqrt{w/\nu}-d_n)=\Phi(-d_n-q\sqrt{w/\nu})\). The integrand of
the first is nonnegative everywhere, and on \(\{w\leq\nu v^2\}\) it is at
least \(\Phi(d_n-qv)\), because \(\sqrt{w/\nu}\leq v\) there and \(\Phi\)
is increasing. Restricting the integral to that set gives
\(\Pr(T>q)\geq\Phi(d_n-qv)G_\nu(\nu v^2)\). At \(\delta=0\),
\(T=Z_c/\sqrt{W/\nu}\) has law \(F_\nu\) by \cref{thm:p15-student-law}, so
\(\Pr(T>t_{\nu,1-\eta})=\eta\) and, by symmetry,
\(\Pr(|T|>t_{\nu,1-\eta/2})=\eta\).
\end{proof}
\end{annexitem}
